\documentclass[11pt,a4paper]{article}

\usepackage[utf8]{inputenc} % allow utf-8 input
\usepackage[T1]{fontenc}    % use 8-bit T1 fonts
\usepackage{hyperref}       % hyperlinks
\usepackage{url}            % simple URL typesetting
\usepackage{booktabs}       % professional-quality tables
\usepackage{amsfonts}       % blackboard math symbols
\usepackage{nicefrac}       % compact symbols for 1/2, etc.
\usepackage{microtype}      % microtypography
\usepackage{xcolor}         % colors
\usepackage{graphicx}       % graphics
\usepackage[normalem]{ulem} % strikethrough

\begin{document}

\textit{\textbf{ Classe de cinquième - Devoir n° 11 - Jeudi 3 décembre 1998 }}

Sujet n° 1

EXERCICE 1 2 points

Compléter :

EXERCICE 2 1 point

Transformer en fraction .

EXERCICE 3 1 point

Compléter :

EXERCICE 4 1 point

Transformer en pourcentage .

EXERCICE 5 3 points

Simplifier les fractions : • • • .

EXERCICE 6 2 points

Comparer les fractions : • •

EXERCICE 7 7 points

Écrire en une seule fraction :

• • •

• • •

EXERCICE 8 2 points

Qu'est-ce qui est le plus long :

“ de mètre ou de mètre ” ? Pourquoi ?

EXERCICE 9 1 point

Vrai ou faux ? : “ 25 a pour moitié 5 ”.

\textit{\textbf{ Classe de cinquième - Devoir n° 11 - Jeudi 3 décembre 1998 }}

Sujet n° 2

EXERCICE 1 2 points

Compléter :

EXERCICE 2 1 point

Transformer en fraction .

EXERCICE 3 1 point

Compléter :

EXERCICE 4 1 point

Transformer en pourcentage .

EXERCICE 5 3 points

Simplifier les fractions : • • • .

EXERCICE 6 2 points

Comparer les fractions : • •

EXERCICE 7 7 points

Écrire en une seule fraction :

• • •

• • •

EXERCICE 8 2 points

Qu'est-ce qui est le plus lourd :

“ de kg ou de kg ” ? Pourquoi ?

EXERCICE 9 1 point

Vrai ou faux ? : “ 4 a pour double 8 ”.

\end{document}